%% ============================================================
%%  week15_merged.tex
%%  Quantum Neural Networks: From Foundations to
%%  Physics-Informed Quantum Circuits
%%
%%  Structure:
%%    0. QNN Recap
%%    1. Input Encoding  (week15.tex)
%%    2. Gradients: Parameter-Shift and Fisher Information
%%    3. Natural Gradient and TDVP
%%    4. Barren Plateaus
%%    5. Classical Physics-Informed Neural Networks
%%    6. Quantum PINNs: Mathematical Foundation
%%    7. Derivative Computation via Shift Rules
%%    8. Poisson Equation: Theory and Code
%%    9. Extended Applications
%%   10. Summary, Exercises, References
%%
%%  Author: Morten Hjorth-Jensen
%%  Department of Physics, University of Oslo
%%  & Michigan State University / FRIB
%%  May 2026
%% ============================================================

\documentclass[aspectratio=169,11pt]{beamer}

%% ── Theme & colours ──────────────────────────────────────────
\usetheme{Madrid}
\usecolortheme{default}
\usefonttheme[onlymath]{serif}

\setbeamercolor{frametitle}{bg=blue!12!white,fg=blue!60!black}
\setbeamercolor{title}{fg=white}
%\setbeamercolor{title}{fg=blue!70!black}
\setbeamercolor{block title}{bg=blue!18!white,fg=blue!60!black}
\setbeamercolor{block body}{bg=blue!4!white}
\setbeamercolor{alerted text}{fg=red!65!black}
\setbeamertemplate{navigation symbols}{}
\setbeamertemplate{footline}{%
  \leavevmode\hbox{%
    \begin{beamercolorbox}[wd=.25\paperwidth,ht=2.25ex,dp=1ex,center]
        {author in head/foot}\usebeamerfont{author in head/foot}
        \insertshortauthor
    \end{beamercolorbox}%
    \begin{beamercolorbox}[wd=.50\paperwidth,ht=2.25ex,dp=1ex,center]
        {title in head/foot}\usebeamerfont{title in head/foot}
        \insertshorttitle
    \end{beamercolorbox}%
    \begin{beamercolorbox}[wd=.25\paperwidth,ht=2.25ex,dp=1ex,right]
        {date in head/foot}\usebeamerfont{date in head/foot}
        \insertshortdate{}\hspace*{1em}%
        \insertframenumber{}/\inserttotalframenumber\hspace*{1ex}
    \end{beamercolorbox}}%
  \vskip0pt}

%% ── Packages ─────────────────────────────────────────────────
\usepackage[utf8]{inputenc}
\usepackage[T1]{fontenc}
\usepackage{amsmath,amssymb,amsfonts,mathtools,bm}
\usepackage{booktabs}
\usepackage{array}
\usepackage{xcolor}
\usepackage{tikz}
\usetikzlibrary{positioning,arrows.meta,shapes.geometric,
                decorations.pathreplacing,fit}
\usepackage{quantikz}
\usepackage{listings}
\usepackage{hyperref}

%% ── Python listing style ─────────────────────────────────────
\lstset{
  language=Python,
  basicstyle=\ttfamily\scriptsize,
  keywordstyle=\color{blue!70!black}\bfseries,
  commentstyle=\color{gray},
  stringstyle=\color{red!60!black},
  breaklines=true,
  frame=single,
  backgroundcolor=\color{blue!3!white},
  xleftmargin=4pt, xrightmargin=4pt,
  numbers=left,
  numberstyle=\ttfamily\tiny\color{gray},
  numbersep=5pt
}

%% ── Macros ────────────────────────────────────────────────────
\newcommand{\R}{\mathbb{R}}
\newcommand{\C}{\mathbb{C}}
\newcommand{\hilb}{\mathcal{H}}
\newcommand{\bx}{\mathbf{x}}
\newcommand{\btheta}{\boldsymbol{\theta}}
\newcommand{\bTheta}{\boldsymbol{\Theta}}
\newcommand{\half}{\tfrac{1}{2}}
\newcommand{\QFI}{\mathbf{F}}
\newcommand{\E}{\mathbb{E}}
\newcommand{\Variance}{\mathrm{Var}}
\newcommand{\manifold}{\mathcal{M}}
\newcommand{\id}{\mathbb{1}}
\newcommand{\Tr}{\operatorname{Tr}}
% Observable expectation shorthand
\newcommand{\evu}[1]{\langle\,#1\,\rangle}

%% ── Metadata ──────────────────────────────────────────────────
\title[QNNs and QPINNs]{Quantum Neural Networks and\\
       Quantum Physics-Informed Neural Networks}
\subtitle{From Input Encoding to PDE Solvers}
\author[M.~Hjorth-Jensen]{Morten Hjorth-Jensen}
\institute[UiO/MSU]{Department of Physics, University of Oslo}
\date{May 6, 2026}

%% ============================================================
\begin{document}
%% ============================================================

\begin{frame}[plain]
  \titlepage
\end{frame}

\begin{frame}[allowframebreaks]{Outline}
  \scriptsize\tableofcontents[hidesubsections]
\end{frame}

%% ============================================================
\section{Recap: Quantum Neural Networks}
%% ============================================================

%% ─── QNN as variational quantum state ───────────────────────
\begin{frame}{QNN as a Variational Quantum State}
  \begin{columns}[T]
    \column{0.52\textwidth}
      \textbf{Quantum Neural Network (QNN)}:
      \begin{align*}
        |\psi(x,\theta)\rangle
        &= U(\theta)\,U(x)\,|0\rangle^{\otimes n}\\
        f(x,\theta)
        &= \langle\psi|\,O\,|\psi\rangle
      \end{align*}
      \begin{itemize}
        \item $U(x)$: \textbf{feature map} --- encodes classical input
        \item $U(\theta)$: \textbf{variational ansatz} --- trainable
        \item $O$: \textbf{observable} --- defines the output scalar
      \end{itemize}
      \begin{block}{Pauli rotation ansatz}
        $U(\theta) = \prod_\ell e^{-i\theta_\ell P_\ell/2}$,
        $P_\ell\in\{X,Y,Z\}^{\otimes n}$
      \end{block}
    \column{0.44\textwidth}
      \small
      \begin{center}
      \begin{quantikz}[row sep=0.26cm, column sep=0.18cm]
        \lstick{$|0\rangle$}
          & \gate{R_x(x_1)} & \gate{R_y(\theta_1)} & \ctrl{1}
          & \gate{R_z(\theta_3)} & \meter{} \\
        \lstick{$|0\rangle$}
          & \gate{R_x(x_2)} & \gate{R_y(\theta_2)} & \targ{}
          & \gate{R_z(\theta_4)} & \meter{}
      \end{quantikz}
      \end{center}
      \begin{alertblock}{Parameter-shift rule}
        $\partial_{\theta_j}f
          = \tfrac{1}{2}[f(\theta_j{+}\tfrac{\pi}{2})
                        -f(\theta_j{-}\tfrac{\pi}{2})]$.
        Exact; 2 evals per param.
      \end{alertblock}
  \end{columns}
\end{frame}

%% ─── QFI / TDVP / Barren Plateaus recap ─────────────────────
\begin{frame}{QFI, TDVP, and Barren Plateaus: Quick Recap}
  \begin{columns}[T]
    \column{0.52\textwidth}
      \textbf{Quantum Fisher Information (QFI):}
      \[
        F_{jk} = 4\,\mathrm{Re}\!\bigl(
          \langle\partial_j\psi|\partial_k\psi\rangle
          - \langle\partial_j\psi|\psi\rangle
            \langle\psi|\partial_k\psi\rangle
        \bigr)
      \]
      Fubini--Study metric on the variational manifold.

      \smallskip
      \textbf{TDVP:} $\sum_k F_{jk}\dot\theta_k
        = \mathrm{Im}\langle\partial_j\psi|H|\psi\rangle$

      \textbf{Natural gradient:}
      $\dot{\btheta} = -\QFI^{-1}\nabla_\theta\mathcal{L}$
    \column{0.44\textwidth}
      \begin{block}{BCH and depth--expressivity}
        \small
        $U^\dagger O U = O + i[H_\theta,O]
          + \tfrac{(i)^2}{2!}[H_\theta,[H_\theta,O]] + \cdots$\\
        Depth $\leftrightarrow$ hierarchy of correlations.
      \end{block}
      \smallskip
      \begin{alertblock}{Barren plateaus}
        Unitary 2-designs give
        $\Variance\!\bigl(\partial_{\theta_j}\mathcal{L}\bigr)
          \sim 1/2^n \to 0$.
        Mitigation: physics-inspired \emph{shallow} ans\"atze
        --- precisely what QPINNs exploit.
      \end{alertblock}
  \end{columns}
\end{frame}

%% ─── Transition frame ─────────────────────────────────────────
\begin{frame}{From Discriminative QNNs to Physics-Informed QNNs}
  \begin{columns}[T]
    \column{0.50\textwidth}
      \begin{block}{Discriminative QNN}
        \begin{itemize}
          \item Input: labelled data pair $(x_k, y_k)$
          \item Loss: $\mathcal{L}=\sum_k(f(x_k;\theta)-y_k)^2$
          \item Learns a mapping from observations
          \item No physical constraint enforced
        \end{itemize}
      \end{block}
      \smallskip
      \begin{alertblock}{Physics-informed QNN (QPINN)}
        \begin{itemize}
          \item Input: collocation points $x_k\in\Omega$
          \item Loss: $\mathcal{L} = \|\mathcal{D}[f_\theta]\|^2
                           +\lambda\|\mathcal{B}[f_\theta]\|^2$
          \item Learns the \emph{solution} of the PDE
          \item Physics is embedded in the loss function
        \end{itemize}
      \end{alertblock}
    \column{0.46\textwidth}
      \textbf{Why quantum circuits for PDEs?}
      \begin{enumerate}\itemsep2pt\small
        \item Exponential compression via entanglement
        \item Natural wave-like encoding via rotation gates
        \item Input-shift rule: spatial derivatives at zero
              extra circuit depth
        \item Connection to quantum simulation of
              differential operators
        \item QFI/TDVP: optimal optimisation geometry
      \end{enumerate}
  \end{columns}
\end{frame}

%% ============================================================
\section{Input Encoding}
%% ============================================================

%% ─── Why encoding matters ────────────────────────────────────
\begin{frame}{Why Encoding Matters}
  A variational quantum circuit takes classical data $\bx\in\R^d$
  and produces a quantum state.  The \textbf{feature map}
  \[
    \Phi:\R^d \longrightarrow \hilb_n = (\C^2)^{\otimes n},
    \qquad \bx \mapsto |\phi(\bx)\rangle,
  \]
  determines:
  \begin{itemize}
    \item what functions the circuit can represent (expressivity),
    \item the induced kernel on the data space,
    \item trainability (gradient norms depend on the map).
  \end{itemize}

  \medskip
  \begin{block}{Principle}
    The encoding is part of the model.  Choosing it well is as
    important as designing the variational ansatz.
    For PDEs, the encoding also determines which spatial frequencies
    the circuit can resolve.
  \end{block}
\end{frame}

%% ─── Basis encoding ──────────────────────────────────────────
\begin{frame}{Basis Encoding}
  \textbf{Idea.} Map a binary string $\bx\in\{0,1\}^n$ to the
  computational basis state:
  \[
    \bx = (x_1,\ldots,x_n)
    \longmapsto |x_1 x_2 \cdots x_n\rangle.
  \]

  \medskip
  \begin{columns}[T]
  \column{0.48\textwidth}
    \textbf{Procedure.}
    \begin{enumerate}
      \item Start in $|0\rangle^{\otimes n}$.
      \item Apply $X_i$ (NOT gate) on qubit $i$ wherever $x_i=1$.
    \end{enumerate}
  \column{0.48\textwidth}
    \textbf{Properties.}
    \begin{itemize}
      \item Only $n$ gates; depth $O(1)$.
      \item Requires $n$ qubits for $n$ bits.
      \item Superposition: encode $N$ strings using
            $\lceil\log_2 N\rceil$ qubits \emph{if} state
            preparation is efficient.
      \item Not suited to continuous features or PDEs.
    \end{itemize}
  \end{columns}
\end{frame}

%% ─── Angle encoding ──────────────────────────────────────────
\begin{frame}{Angle Encoding}
  \textbf{Idea.} Encode $x_i\in\R$ into the rotation angle of qubit $i$:
  \[
    |\phi(\bx)\rangle
    = \bigotimes_{i=1}^{n} R_y(\pi x_i)|0\rangle
    = \bigotimes_{i=1}^{n}
      \begin{pmatrix}\cos(\pi x_i/2)\\ \sin(\pi x_i/2)\end{pmatrix}.
  \]

  \medskip
  Rotation matrices for $k\in\{x,y,z\}$:
  $R_k(\alpha) = e^{-i\alpha\sigma_k/2}$,
  \[
    R_y(\alpha)=\begin{pmatrix}\cos\frac\alpha2 & -\sin\frac\alpha2\\
                               \sin\frac\alpha2 &  \cos\frac\alpha2
               \end{pmatrix},
    \quad
    R_z(\alpha)=\begin{pmatrix}e^{-i\alpha/2}&0\\0&e^{i\alpha/2}
               \end{pmatrix}.
  \]

  \medskip
  \begin{itemize}
    \item One feature per qubit; $n=d$.  Depth $O(1)$.
    \item Product state --- \emph{no entanglement yet}; entanglement
          added by two-qubit gates in the variational layers.
    \item The output $f(x;\theta)$ is a \emph{Fourier series} in $x$
          with frequencies determined by the generator spectrum.
  \end{itemize}
\end{frame}

%% ─── Amplitude encoding ──────────────────────────────────────
\begin{frame}{Amplitude Encoding}
  \textbf{Idea.} Embed $\bx\in\R^{2^n}$, $\|\bx\|=1$, into
  state amplitudes:
  \[
    |\phi(\bx)\rangle = \sum_{j=0}^{2^n-1} x_j |j\rangle.
  \]

  \begin{columns}[T]
  \column{0.52\textwidth}
    \textbf{Capacity.}
    \begin{itemize}
      \item $n$ qubits encode $2^n$ features ---
            \emph{exponential compression}.
      \item Exploits the full Hilbert space.
    \end{itemize}
    \medskip
    \textbf{Challenges.}
    \begin{itemize}
      \item General state preparation requires $O(2^n)$ gates.
      \item Normalisation: $\sum_j x_j^2=1$.
      \item $N$ data points need $N$ circuit runs.
    \end{itemize}
  \column{0.44\textwidth}
    \begin{alertblock}{Read-out bottleneck}
      The exponential compression does \emph{not} automatically
      give a speedup: reading out all $2^n$ amplitudes requires
      exponentially many measurements.
      This is the fundamental quantum read-out bottleneck
      relevant to all QPINNs.
    \end{alertblock}
  \end{columns}
\end{frame}

%% ─── IQP and ZZ feature maps ─────────────────────────────────
\begin{frame}{Structured Feature Maps: IQP and ZZ}
  \textbf{IQP feature map} (Havl\'{\i}\v{c}ek et al., Nature 2019):
  \[
    |\phi(\bx)\rangle = U_\Phi(\bx)|+\rangle^{\otimes n},
    \quad
    U_\Phi(\bx) = e^{i\phi(\bx)},
    \quad
    \phi(\bx) = \sum_S c_S(\bx)\bigotimes_{i\in S}Z_i.
  \]

  \medskip
  \textbf{Second-order ZZ feature map} (two layers):
  \begin{align*}
    U_\Phi^{(2)}(\bx)
    &= V_2(\bx)H^{\otimes n}V_1(\bx)H^{\otimes n},\\
    V_k(\bx)
    &= \exp\!\Bigl(i\sum_i x_i Z_i
       + i\sum_{i<j}(\pi-x_i)(\pi-x_j)Z_iZ_j\Bigr).
  \end{align*}

  \medskip
  These encode two-body correlations; the resulting kernel
  $K(\bx,\bx')=\bigl|\langle\phi(\bx)|\phi(\bx')\rangle\bigr|^2$
  is believed to be classically hard to simulate.
  For QPINNs, IQP encoding introduces cross-frequency terms
  $\omega_i\omega_j$ that enrich the Fourier spectrum of
  the solution approximation.
\end{frame}

%% ─── Data Re-uploading ───────────────────────────────────────
\begin{frame}{Data Re-Uploading: Expressivity via Repetition}
\small
  \textbf{Key idea} (P\'erez-Salinas et al.\ 2020): encode
  $\bx$ \emph{at every layer}:
  \[
    U(\bx,\bTheta) = \prod_{\ell=1}^{L}
      \Bigl[W(\bTheta^{(\ell)})\,S(\bx)\Bigr],
    \quad S(\bx) = \bigotimes_i R_y(\pi x_i).
  \]

  \begin{columns}[T]
  \column{0.52\textwidth}
    \textbf{Why it helps.}
    \begin{itemize}\itemsep1pt\small
      \item $L$ re-uploads $\Rightarrow$ Fourier series up to freq.\ $L$.
      \item For PDEs: $\omega_{\max} = Ln/2$ ($n$ qubits, $L$ layers).
    \end{itemize}
    \textbf{Fourier:} $f(\bx;\bTheta)
      = \sum_{\bm\omega\in\Omega}c_{\bm\omega}(\bTheta)e^{i\bm\omega\cdot\bx}$.
    Frequencies fixed by circuit; coefficients by $\bTheta$.
  \column{0.44\textwidth}
    \small
    \begin{block}{Encoding comparison for PDEs}
      \begin{center}
      \renewcommand{\arraystretch}{1.1}
      \begin{tabular}{lcc}
        \toprule
        \textbf{Encoding} & \textbf{Qubits} & $\omega_{\max}$ \\
        \midrule
        Angle (1 layer) & $n$   & $n/2$ \\
        Re-uploading    & $n$   & $Ln/2$ \\
        IQP / ZZ        & $n$   & $n^2/4$ \\
        Amplitude       & $n$   & $2^n/2$ \\
        \bottomrule
      \end{tabular}
      \end{center}
    \end{block}
    \begin{alertblock}{Design rule}
      Choose $n$, $L$ s.t.\ $Ln \geq 2\omega^\star$.
    \end{alertblock}
  \end{columns}
\end{frame}

%% ============================================================
\section{Gradient Computation: Parameter-Shift and Fisher Information}
%% ============================================================

%% ─── VQC setup ───────────────────────────────────────────────
\begin{frame}{Variational Quantum Circuits: Setup}
  For a parameterised circuit
  $U(\btheta) = \prod_j G_j(\theta_j)$ with
  $G_j(\theta_j) = e^{-i\theta_j P_j/2}$, $P_j^2=\id$:
  \[
    f(\btheta)
    = \langle\psi_0|U^\dagger(\btheta)\,\hat{B}\,
      U(\btheta)|\psi_0\rangle.
  \]

  \begin{block}{Parameter-shift rule: derivation}
    Write $U(\theta_j)=A_j G_j(\theta_j)B_j$,
    $G_j=\cos\tfrac{\theta_j}{2}\id-i\sin\tfrac{\theta_j}{2}P_j$:
    \[
      \frac{\partial f}{\partial\theta_j}
      = \frac{1}{2}\bigl[f(\theta_j+\tfrac{\pi}{2})
                        -f(\theta_j-\tfrac{\pi}{2})\bigr].
    \]
    Exact; 2 circuit evaluations per parameter.
    Generalises to $r[f(\theta_j+\frac\pi{4r})-f(\theta_j-\frac\pi{4r})]$
    for eigenvalues $\pm r$.
  \end{block}
  \vspace{-2pt}
  \begin{columns}[T]
  \column{0.5\textwidth}
    \textbf{2nd deriv.:}
    $\partial^2_{\theta_j}f
      = f(\theta_j{+}\pi) - 2f(\theta_j) + f(\theta_j{-}\pi)$.
  \column{0.46\textwidth}
    Used in QFI and Hessian optimisers;
    combined with input-shift gives mixed PDE derivatives.
  \end{columns}
\end{frame}

%% ─── QFI definition ──────────────────────────────────────────
\begin{frame}{Quantum Fisher Information Matrix}
\small
  \textbf{Definition (pure state $|\psi(\btheta)\rangle$):}
  \[
    F_{jk}(\btheta)
    = 4\,\mathrm{Re}\!\bigl(
        \langle\partial_j\psi|\partial_k\psi\rangle
       -\langle\partial_j\psi|\psi\rangle
        \langle\psi|\partial_k\psi\rangle
      \bigr)
  \]
  where $|\partial_j\psi\rangle=\partial|\psi(\btheta)\rangle/\partial\theta_j$.

  \smallskip
  \begin{columns}[T]
  \column{0.52\textwidth}
    \small\textbf{Geometric meaning.}  $\QFI$ is the
    \textbf{Fubini--Study metric} on the variational manifold:
    \[
      ds^2 = \sum_{jk}F_{jk}\,d\theta_j\,d\theta_k.
    \]
    It measures the \emph{information-geometric distance}
    between neighbouring quantum states.

    \textbf{Covariance form.} For product-of-rotations circuits:
    \[
      F_{jk} = 4\,\mathrm{Cov}_\psi(G_j, G_k).
    \]
  \column{0.44\textwidth}
    \begin{block}{Circuit formula (2-fold shifts)}
      For overlap $\mathcal{O}(\btheta)=|\langle\psi_{\rm ref}|\psi\rangle|^2$:
      \begin{align*}
        F_{jk} &= \tfrac{1}{2}[\mathcal{O}_{j+,k+}
         -\mathcal{O}_{j+,k-}\\
         &\quad-\mathcal{O}_{j-,k+}
         +\mathcal{O}_{j-,k-}].
      \end{align*}
      Shift each index by $\pm\tfrac\pi2$.
      Cost: $O(p^2)$ circuit evals.
    \end{block}
  \end{columns}
\end{frame}

%% ─── Natural gradient and TDVP ───────────────────────────────
\begin{frame}{Natural Gradient and the TDVP}
  \begin{columns}[T]
  \column{0.52\textwidth}
    \textbf{Natural gradient} (Amari 1998; Stokes et al.\ 2020):
    \[
      \btheta \leftarrow \btheta
        - \eta\,\QFI(\btheta)^{+}\,\nabla_\theta C.
    \]
    \textbf{Derivation.}  Minimise $C(\btheta+\delta\btheta)$
    subject to fixed Bures distance:
    $\delta\btheta^\top\QFI\,\delta\btheta = \varepsilon^2$,
    giving $\delta\btheta\propto-\QFI^{-1}\nabla C$.

    The update is \emph{reparameterisation-invariant}: it always
    descends the steepest direction on $\manifold$.

  \column{0.44\textwidth}
    \begin{block}{TDVP equations of motion}
      From the Dirac--Frenkel principle, projecting
      $(i\partial_t - H)|\psi\rangle = 0$
      onto the tangent space $\{|\partial_j\psi\rangle\}$:
      \[
        \sum_k F_{jk}\,\dot{\theta}_k
        = 2\,\mathrm{Im}\langle\partial_j\psi|H|\psi\rangle
      \]
    \end{block}
    \vspace{0.15cm}
    \begin{alertblock}{Imaginary-time TDVP $=$ natural gradient VQE}
      Substituting $t\to-i\tau$:
      $\dot{\btheta} = -\QFI^{-1}\nabla_\theta E(\btheta)$.
      Equivalent to Stochastic Reconfiguration (SR) in VMC.
    \end{alertblock}
  \end{columns}
\end{frame}

%% ─── Effective dimension ─────────────────────────────────────
\begin{frame}{Effective Dimension and Expressivity}
  Abbas et al.\ (2021) define the \textbf{effective dimension}
  via the average log-volume of the QFI:
  \[
    d_{\rm eff}
    = 2\,\frac{1}{\log N}
      \log\!\left[\frac{1}{V}\int_{\btheta}
        \sqrt{\det\!\left(
          I + \frac{N}{2\pi\ln N}\,\QFI(\btheta)
        \right)}\,d\btheta\right],
  \]
  where $N$ is the number of training samples.

  \medskip
  \begin{columns}[T]
  \column{0.52\textwidth}
    \begin{itemize}
      \item High $d_{\rm eff}$ $\Rightarrow$ model fits many
            functions; good generalisation capacity.
      \item Some QNNs achieve higher $d_{\rm eff}$ than
            classical networks of comparable parameter count.
      \item $\QFI\approx 0$ (rank-deficient) $\Rightarrow$
            parameters are redundant; barren-plateau-like
            flat directions.
    \end{itemize}
  \column{0.44\textwidth}
    \begin{block}{QFI as trainability indicator}
      \[
        \left|\frac{\partial C}{\partial\theta_j}\right|^2
        \leq \lambda_{\max}(\QFI)\cdot\Variance_\psi(\hat{B})
      \]
      (Cauchy--Schwarz on the Fubini--Study metric).\\[4pt]
      Small $\lambda_{\max}(\QFI)$
      $\Leftrightarrow$ all gradients are small
      $\Leftrightarrow$ barren plateau.
    \end{block}
  \end{columns}
\end{frame}

%% ============================================================
\section{Barren Plateaus}
%% ============================================================

%% ─── BP definition ───────────────────────────────────────────
\begin{frame}{Barren Plateaus: Definition and Origin}
\small
  \begin{block}{Definition (McClean et al.\ 2018)}
    $\{U(\btheta)\}$ has a \textbf{barren plateau} if:
    \[
      \E_{\btheta}\!\left[\partial_{\theta_j}C\right] = 0,
      \qquad
      \Variance_{\btheta}\!\left[\partial_{\theta_j}C\right]
      = O(b^{-n}),\; b>1.
    \]
  \end{block}

  \begin{columns}[T]
  \column{0.52\textwidth}
    \textbf{Mathematical origin.}
    For Haar-random $U$ on $\mathrm{SU}(2^n)$:
    \[
      \E_U[\langle\hat{O}\rangle]
      = \tfrac{\Tr\hat{O}}{2^n},\quad
      \Variance[\partial_{\theta_j}C]
      \sim \tfrac{\Tr\hat{O}^2}{4^n}.
    \]
    For a local observable $\hat{O}=Z_k$:
    $\Variance \sim 1/2^n$.

  \column{0.44\textwidth}
    \textbf{Concentration of measure} (L\'evy's lemma):
    in high dimensions, all Lipschitz functions
    concentrate near their mean with probability
    $\geq 1 - 2e^{-\Omega(2^{2n}\varepsilon^2)}$.

    \begin{alertblock}{Consequence}
      Deep random circuits make $C(\btheta)\approx\Tr\hat{O}/2^n$
      almost everywhere: landscape is flat.
      This is \emph{not} a local-minima problem.
    \end{alertblock}
  \end{columns}
\end{frame}

%% ─── Noise-induced BP ────────────────────────────────────────
\begin{frame}{Noise-Induced Barren Plateaus and Mitigation}
  \begin{columns}[T]
  \column{0.52\textwidth}
    \textbf{Hardware noise} (Wang et al.\ 2021): depolarising
    noise with rate $\epsilon$ per gate gives:
    \[
      |\langle\hat{O}\rangle_{\rm noisy}|
      \leq (1-\epsilon)^{pd}\,
      |\langle\hat{O}\rangle_{\rm ideal}|,
    \]
    and $\Variance[\partial_{\theta_j}C_{\rm noisy}]
      \sim (1-\epsilon)^{2pd}\Variance[\partial_{\theta_j}C_{\rm ideal}]$.
    Affects even \emph{shallow} circuits with many gates.
  \column{0.44\textwidth}
    \begin{block}{Mitigation strategies}
      \footnotesize
      \begin{tabular}{ll}
        \textbf{Strategy} & \textbf{Mechanism} \\
        \midrule
        Local cost & poly$(n)$ scaling \\
        Problem ansatz & exploit symmetries \\
        Layerwise training & $O(1)$ depth locally \\
        Identity init & $O(1)$ init gradients \\
        Natural gradient & QFI-respecting steps \\
        Shallow circuits & avoid Haar regime \\
      \end{tabular}
    \end{block}
    \begin{alertblock}{QPINNs and BPs}
      Physics-inspired shallow ans\"atze naturally avoid
      the random circuit regime --- PDE structure
      guides the ansatz design.
    \end{alertblock}
  \end{columns}
\end{frame}

%% ─── Entanglement and QNTK ───────────────────────────────────
\begin{frame}{Entanglement, QNTK, and Training Regimes}
  \begin{columns}[T]
  \column{0.52\textwidth}
    \textbf{Entanglement entropy:}
    $S(\rho_A) = -\Tr(\rho_A\log\rho_A)$
    governs the training regime:

    \begin{center}
    \begin{tabular}{lll}
      \toprule
      Regime & $S(\rho_A)$ & Training \\
      \midrule
      Low & $\approx0$ & Classically simulable \\
      Volume-law & $\Theta(n)$ & Barren plateaus \\
      Area-law & $O(|\partial A|)$ & Trainable \\
      \bottomrule
    \end{tabular}
    \end{center}

  \column{0.44\textwidth}
    \textbf{Quantum Neural Tangent Kernel (QNTK):}
    \[
      \Theta(\bx,\bx')
      = \sum_j
        \frac{\partial f(\bx;\btheta)}{\partial\theta_j}
        \frac{\partial f(\bx';\btheta)}{\partial\theta_j}.
    \]
    \begin{itemize}
      \item Governs training dynamics in the linearised
            regime.
      \item At a barren plateau:
            $\Theta\approx 0$ for all $\bx,\bx'$.
      \item Bound: $\Theta(\bx,\bx)
            \leq\lambda_{\max}(\QFI)$.
      \item Connects QNN training to quantum kernel
            methods (QSVMs).
    \end{itemize}
  \end{columns}
\end{frame}

%% ============================================================
\section{Classical Physics-Informed Neural Networks}
%% ============================================================

%% ─── Classical PINN ──────────────────────────────────────────
\begin{frame}{Classical PINNs: Raissi, Perdikaris, Karniadakis (2019)}
\small
  \begin{columns}[T]
  \column{0.52\textwidth}
    A classical PINN approximates the solution of
    $\mathcal{N}[u](x,t)=0$ on $\Omega$ by a neural network
    $u_\theta(x,t)$ minimising:
    \begin{align*}
      \mathcal{L}(\theta)
      &= \underbrace{\tfrac{1}{N_r}\textstyle\sum_k
          |\mathcal{N}[u_\theta](x_k,t_k)|^2}_{\mathcal{L}_r}\\
      &\;+ \lambda\underbrace{\tfrac{1}{N_b}\textstyle\sum_k
          |u_\theta(x_k^b,t_k^b)-g_k|^2}_{\mathcal{L}_b}
    \end{align*}
    Collocation points $\{x_k\}$ sampled in $\Omega$;
    $\{x_k^b\}$ on $\partial\Omega$.
  \column{0.44\textwidth}
    \begin{block}{Automatic differentiation}\footnotesize
      Classical PINNs use autograd (no finite differences).
    \end{block}
    \begin{alertblock}{Spectral bias}
      Classical NNs learn low-frequency components much
      faster than high-frequency ones ---
      motivation for QPINNs with explicit frequency spectra.
    \end{alertblock}
    \begin{block}{Examples}\footnotesize
      Poisson $-\nabla^2 u\!=\!f$;\;
      Heat $\partial_t u\!=\!\kappa\nabla^2 u$;\;
      Wave $\partial_{tt}u\!=\!c^2\nabla^2 u$;\;
      Burgers $\partial_t u\!+\!u\partial_x u\!=\!\nu\partial_{xx}u$;\;
      Schr\"odinger $i\partial_t\psi\!=\!H\psi$.
    \end{block}
  \end{columns}
\end{frame}

%% ─── PINN loss landscape ────────────────────────────────────
\begin{frame}{PINN Loss Landscape and the Neural Tangent Kernel}
  \begin{columns}[T]
  \column{0.52\textwidth}
    Residual at collocation point $x_k$:
    $r_k(\theta) = \mathcal{D}[f_\theta](x_k)$.
    \[
      \mathcal{L}(\theta)
      = \frac{1}{N}\sum_k r_k^2,\quad
      \nabla_\theta\mathcal{L}
      = \frac{2}{N}\sum_k r_k\,\nabla_\theta r_k.
    \]
    The \textbf{Neural Tangent Kernel (NTK)} governs
    training dynamics:
    \[
      \Theta_{kl}
      = \nabla_\theta f_\theta(x_k)^\top
        \nabla_\theta f_\theta(x_l).
    \]
    Eigenvalues of $\Theta$ determine convergence rates
    per frequency component.
  \column{0.44\textwidth}
    \begin{alertblock}{Failure modes}
      \begin{enumerate}
        \item \textbf{Spectral bias:} high-frequency modes
              converge slowly
        \item \textbf{Loss balance:}
              $\mathcal{L}_r$ and $\mathcal{L}_b$ compete
        \item \textbf{Stiff equations:} chaotic landscape
        \item \textbf{High dimensionality:} exponential cost
      \end{enumerate}
    \end{alertblock}
    \begin{block}{Quantum motivation}
      QPINNs address points 1 and 4: richer frequency
      spectra via entanglement; quantum memory scales as
      $O(2^n)$ with $n$ qubits.
    \end{block}
  \end{columns}
\end{frame}

%% ============================================================
\section{Quantum PINNs: Mathematical Foundation}
%% ============================================================

%% ─── QPINN ansatz ────────────────────────────────────────────
\begin{frame}{QPINN: The Quantum Ansatz for PDE Solutions}
  \begin{columns}[T]
  \column{0.52\textwidth}
    Replace the classical network $u_\theta(x)$ with a
    variational quantum circuit output:
    \[
      u_\theta(x)
      = \langle 0|U^\dagger(x,\theta)\,O\,U(x,\theta)|0\rangle
    \]
    where:
    \begin{itemize}
      \item $U(x,\theta) = W_L(\theta)\cdots W_1(\theta)\,S(x)$
      \item $S(x) = \bigotimes_i R_y(\pi x)$: data encoding
      \item $W_l(\theta)$: variational Pauli layers
      \item $O = Z_0$: fixed observable giving scalar output
    \end{itemize}
    The output $u_\theta(x)\in[-1,1]$ approximates the
    normalised PDE solution.
  \column{0.44\textwidth}
    \begin{center}
    \begin{quantikz}[row sep=0.30cm, column sep=0.28cm]
      \lstick{$|0\rangle$}
        & \gate{S_1} & \gate{W_1} & \ctrl{1}
        & \gate{W_1'} & \meter{$O$} \\
      \lstick{$|0\rangle$}
        & \gate{S_2} & \gate{W_2} & \targ{}
        & \gate{W_2'} & \qw \\
      \lstick{$|0\rangle$}
        & \gate{S_3} & \gate{W_3} & \ctrl{1}
        & \gate{W_3'} & \qw \\
      \lstick{$|0\rangle$}
        & \gate{S_4} & \gate{W_4} & \targ{}
        & \gate{W_4'} & \qw
    \end{quantikz}
    \end{center}
    {\small $S_i$: encoding; $W_i$: variational; $O=Z_0$.}
  \end{columns}
\end{frame}

%% ─── QPINN loss ──────────────────────────────────────────────
\begin{frame}{The QPINN Loss Function}
  For $\mathcal{D}[u](x) = f(x)$ on $\Omega$ with
  $u=g$ on $\partial\Omega$:
  \[
    \mathcal{L}(\theta)
    = \underbrace{\frac{1}{N_r}\sum_k
        \bigl|\mathcal{D}[u_\theta](x_k) - f(x_k)\bigr|^2}_{
        \mathcal{L}_r}
      + \lambda\,\underbrace{\frac{1}{N_b}\sum_k
        \bigl|u_\theta(x_k^b) - g(x_k^b)\bigr|^2}_{
        \mathcal{L}_b}.
  \]

  \begin{columns}[T]
  \column{0.52\textwidth}
    \textbf{Gradient w.r.t.\ $\theta_\mu$:}
    \[
      \frac{\partial\mathcal{L}_r}{\partial\theta_\mu}
      = \frac{2}{N_r}\sum_k r_k\,
        \frac{\partial\mathcal{D}[u_\theta](x_k)}
             {\partial\theta_\mu}.
    \]
    For a \emph{linear} operator $\mathcal{D}$:
    \[
      \frac{\partial\mathcal{D}[u_\theta]}{\partial\theta_\mu}
      = \mathcal{D}\!\left[
          \frac{\partial u_\theta}{\partial\theta_\mu}
        \right],
    \]
    so $\mathcal{D}$ and the parameter-shift commute.
  \column{0.44\textwidth}
    \begin{alertblock}{Two shift rules working together}
      \begin{itemize}
        \item \textbf{Input-shift}: computes spatial
              derivatives $\mathcal{D}[u_\theta]$
        \item \textbf{Parameter-shift}: computes gradients
              $\partial u_\theta/\partial\theta_\mu$
      \end{itemize}
      Total cost per gradient step:
      $O(N_p \times N_r)$ circuit evaluations.
    \end{alertblock}
  \end{columns}
\end{frame}

%% ============================================================
\section{Derivative Computation via Shift Rules}
%% ============================================================

%% ─── Input-shift rule ────────────────────────────────────────
\begin{frame}{The Input-Shift Rule: Spatial Derivatives}
  For angle encoding $S(x) = e^{-ixP/2}$ (Pauli $P$,
  eigenvalues $\pm\frac{1}{2}$):

  \begin{columns}[T]
  \column{0.52\textwidth}
    \textbf{First derivative} (input-shift rule):
    \[
      \frac{\partial u_\theta}{\partial x}
      = \tfrac{1}{2}\bigl[
          u_\theta(x{+}\tfrac{\pi}{2})
         -u_\theta(x{-}\tfrac{\pi}{2})
        \bigr]
    \]
    \textbf{Second derivative:}
    \[
      \frac{\partial^2 u_\theta}{\partial x^2}
      = u_\theta(x{+}\tfrac{\pi}{2})
        - 2u_\theta(x)
        + u_\theta(x{-}\tfrac{\pi}{2})
    \]
    \textbf{$n$-th order:}
    \[
      \partial_x^n u_\theta
      = \sum_{k=0}^n(-1)^k\binom{n}{k}
        u_\theta\!\left(x+\tfrac{(n-2k)\pi}{4}\right)
    \]
  \column{0.44\textwidth}
    \begin{alertblock}{Key result}
      All spatial derivatives of $u_\theta(x)$ are computable
      by evaluating the \emph{same circuit} at shifted input
      values.  No ancilla qubits; no additional circuit depth.
    \end{alertblock}
    \vspace{0.15cm}
    \begin{block}{Cost budget}
      $k$-th order derivative at one point: $k+1$ circuit runs.
      \\[4pt]
      Laplacian $\partial_{xx}$: \textbf{3 runs} at
      $x{-}\tfrac\pi2$, $x$, $x{+}\tfrac\pi2$.
    \end{block}
  \end{columns}
\end{frame}

%% ─── Unified shift framework ─────────────────────────────────
\begin{frame}{Unified Shift Framework: One Rule, Two Applications}
  For \emph{any} gate $e^{-i\phi P/2}$ (Pauli $P$):
  \[
    \frac{\partial}{\partial\phi}\langle O\rangle_\phi
    = \frac{1}{2}\bigl[
        \langle O\rangle_{\phi+\pi/2}
       -\langle O\rangle_{\phi-\pi/2}
      \bigr].
  \]

  \begin{columns}[T]
  \column{0.52\textwidth}
    \textbf{Two interpretations:}
    \begin{itemize}
      \item $\phi = \theta_\mu$ (variational parameter):
            \textbf{parameter-shift} --- gradient for training
      \item $\phi = x$ (input encoding):
            \textbf{input-shift} --- spatial derivative of PDE solution
    \end{itemize}
    They are \emph{the same formula}, applied to different gates
    in the circuit.

    \begin{center}
    \begin{quantikz}[row sep=0.32cm, column sep=0.30cm]
      \lstick{$|0\rangle$}
        & \gate{R_y(x\pm\frac\pi2)} & \gate{W(\theta)}
        & \meter{}
    \end{quantikz}
    \end{center}
    {\small Input-shift: $\partial_x u_\theta$.}
  \column{0.44\textwidth}
    \begin{center}
    \begin{quantikz}[row sep=0.32cm, column sep=0.30cm]
      \lstick{$|0\rangle$}
        & \gate{R_y(x)} & \gate{R_y(\theta\pm\frac\pi2)}
        & \meter{}
    \end{quantikz}
    \end{center}
    {\small Parameter-shift: $\partial_\theta u_\theta$.}

    \smallskip
    \begin{block}{Commutativity (linear PDEs)}
      The shifts commute:
      $\partial_{\theta_\mu}\mathcal{D}[u_\theta](x)
        = \mathcal{D}[\partial_{\theta_\mu}u_\theta](x)$.
      This is why QPINN training is tractable.
    \end{block}
  \end{columns}
\end{frame}

%% ─── Multi-dimensional PDEs ──────────────────────────────────
\begin{frame}{Multi-Dimensional PDEs and the Variational Form}
  \begin{columns}[T]
  \column{0.52\textwidth}
    \textbf{Multi-dimensional encoding.}
    For $u(x_1,\ldots,x_d)$ with $n_i$ qubits for $x_i$:
    \[
      S(\bx) = \bigotimes_{i=1}^d R_y(\pi x_i)^{\otimes n_i}.
    \]
    Cross-derivatives: shift simultaneously in $x_i$ and $x_j$:
    $\partial_{x_i}\partial_{x_j}u_\theta$ via simultaneous input-shifts.
    Total qubits: $n = n_1+\cdots+n_d$.
    Classical PINN: $O(\varepsilon^{-d})$ parameters;
    QPINN: $O(d)$ qubits if solution has tensor-product
    structure.

  \column{0.44\textwidth}
    \begin{block}{Variational (Ritz) form}
      For Poisson $-\nabla^2 u = f$, minimise:
      \[
        \mathcal{L}_{\rm Ritz}
        = \sum_k w_k\bigl[
            \tfrac12(u_\theta'(x_k))^2
            - f_k u_k
          \bigr],
      \]
      Only first derivatives needed; smoother landscape.
    \end{block}
    \begin{alertblock}{Ritz-QPINN = quantum Galerkin}
      The ansatz $u_\theta(x)$ plays the finite element
      basis role but spans an exponentially large function space.
    \end{alertblock}
  \end{columns}
\end{frame}

%% ============================================================
\section{Poisson Equation: Theory and Code}
%% ============================================================

%% ─── 1D Poisson ──────────────────────────────────────────────
\begin{frame}{The 1D Poisson Equation}
  \begin{columns}[T]
  \column{0.52\textwidth}
    \textbf{Problem:}
    \[
      -\frac{d^2u}{dx^2} = f(x),\quad x\in(0,1),
      \quad u(0)=u(1)=0.
    \]
    Exact solution for $f(x)=\pi^2\sin(\pi x)$:
    \[
      u^\star(x) = \sin(\pi x).
    \]

    \textbf{QPINN residual loss:}
    \[
      \mathcal{L}_r
      = \frac{1}{N_r}\sum_{k=1}^{N_r}
        \left[-\frac{d^2u_\theta}{dx^2}\bigg|_{x_k}
              - f(x_k)\right]^2.
    \]
    \textbf{Boundary loss:}
    $\mathcal{L}_b = u_\theta(0)^2 + u_\theta(1)^2$.
  \column{0.44\textwidth}
    \textbf{Second derivative via input-shift:}
    \[
      \frac{d^2u_\theta}{dx^2}\bigg|_x
      = u_\theta(x{+}\tfrac\pi2)
        - 2u_\theta(x)
        + u_\theta(x{-}\tfrac\pi2).
    \]
    This uses $S(x)=e^{-i\pi xP/2}$ with $P=Y\otimes\cdots$:
    exactly 3 circuit evaluations per collocation point.
    \begin{alertblock}{Circuit count per gradient step}
      $N_r$ interior points $\times\,3$ (Laplacian)
      $\times\,2N_p$ (parameter shift)
      $= 6N_pN_r$ total.
      For $N_p=8$, $N_r=20$: 960 evaluations per step.
    \end{alertblock}
  \end{columns}
\end{frame}

%% ─── Python code: gates and circuit ──────────────────────────
\begin{frame}[fragile]{Python: Circuit and Gate Definitions}
\begin{lstlisting}[basicstyle=\tiny\ttfamily]
import numpy as np

def ry(t):
    c, s = np.cos(t / 2), np.sin(t / 2)
    return np.array([[c, -s], [s, c]], dtype=complex)

I2   = np.eye(2, dtype=complex)
CNOT = np.array([[1,0,0,0],[0,1,0,0],[0,0,0,1],[0,0,1,0]], dtype=complex)
Z0   = np.kron(np.diag([1., -1.]).astype(complex), I2)  # Z on qubit 0

def var_layer(p):
    """Ry(p0)xRy(p1) -> CNOT -> Ry(p2)xRy(p3)"""
    return CNOT @ np.kron(ry(p[0]), ry(p[1]))

def circuit(x, params):
    """
    2-qubit QPINN circuit for 1D PDE.
    Encoding:  Ry(pi*x) on each qubit.
    Ansatz:    2 variational layers (8 parameters).
    Output:    <Z_0> in [-1, 1].
    """
    state = np.zeros(4, dtype=complex); state[0] = 1.0
    U = np.kron(ry(np.pi * x), ry(np.pi * x))       # angle encoding
    U = var_layer(params[:4]) @ U                    # layer 1
    U = var_layer(params[4:8]) @ U                   # layer 2
    psi = U @ state
    return float((psi.conj() @ Z0 @ psi).real)
\end{lstlisting}
\end{frame}

%% ─── Python code: derivatives and loss ──────────────────────
\begin{frame}[fragile]{Python: Input-Shift Derivatives and Loss}
\begin{lstlisting}[basicstyle=\scriptsize\ttfamily]
SHIFT = np.pi / 2

def du_dx(x, params):
    return 0.5 * (circuit(x + SHIFT, params)
                - circuit(x - SHIFT, params))

def d2u_dx2(x, params):
    return (circuit(x + SHIFT, params)
          - 2 * circuit(x, params)
          + circuit(x - SHIFT, params))

x_r, x_b, lambda_bc = np.linspace(0.05,0.95,20), np.array([0.,1.]), 10.

def loss(params):
    res = [-d2u_dx2(x,params) - np.pi**2*np.sin(np.pi*x) for x in x_r]
    L_r = np.mean(np.array(res)**2)
    L_b = circuit(x_b[0],params)**2 + circuit(x_b[1],params)**2
    return L_r + lambda_bc * L_b
\end{lstlisting}
\end{frame}

\begin{frame}[fragile]{Python: Gradient via Parameter-Shift Rule}
\begin{lstlisting}[basicstyle=\scriptsize\ttfamily]
def grad_loss(params, shift=np.pi / 2):
    """Exact gradient via parameter-shift rule."""
    g = np.zeros(len(params))
    for mu in range(len(params)):
        p_plus  = params.copy(); p_plus[mu]  += shift
        p_minus = params.copy(); p_minus[mu] -= shift
        g[mu] = 0.5 * (loss(p_plus) - loss(p_minus))
    return g
\end{lstlisting}
\end{frame}

%% ─── Python code: training ───────────────────────────────────
\begin{frame}[fragile]{Python: Training Loop (Adam Optimiser)}
\begin{lstlisting}[basicstyle=\tiny\ttfamily]
np.random.seed(7)
params = np.random.uniform(0, np.pi, 8)    # 2 layers x 4 params

# Adam hyperparameters
lr, b1, b2, eps = 0.15, 0.9, 0.999, 1e-8
m, v = np.zeros(8), np.zeros(8)
history = []

for step in range(300):
    g = grad_loss(params)
    t = step + 1
    m  = b1 * m + (1 - b1) * g
    v  = b2 * v + (1 - b2) * g**2
    mc = m / (1 - b1**t);  vc = v / (1 - b2**t)
    params -= lr * mc / (np.sqrt(vc) + eps)
    history.append(loss(params))
    if (step + 1) % 75 == 0:
        print(f"Step {step+1:4d}  loss = {history[-1]:.6f}")

# Evaluate and compare
x_test  = np.linspace(0, 1, 100)
u_qpinn = np.array([circuit(x, params) for x in x_test])
u_exact = np.sin(np.pi * x_test)
l2_err  = np.sqrt(np.mean((u_qpinn - u_exact)**2))
print(f"\nL2 error vs. sin(pi*x): {l2_err:.4f}")
# Expected: L2 error ~ 0.03 -- 0.05 after 300 Adam steps
\end{lstlisting}
\end{frame}

%% ─── Results and interpretation ────────────────────────────
\begin{frame}{Results and Interpretation}
  \begin{columns}[T]
  \column{0.52\textwidth}
    \textbf{Expected output:}
    \begin{itemize}\itemsep1pt
      \item L2 error $\lesssim 0.05$ after 300 steps
      \item Loss drops 2--3 orders of magnitude
      \item QPINN traces $\sin(\pi x)$ accurately
    \end{itemize}
    \textbf{Improving accuracy:}
    \begin{itemize}\itemsep1pt
      \item More qubits (wider frequency spectrum)
      \item More variational layers (richer coefficients)
      \item More collocation points $N_r$
      \item Natural gradient (QFI) instead of Adam
    \end{itemize}

  \column{0.44\textwidth}
    \begin{block}{Connection to Section 2}
      Adam uses Euclidean geometry.
      Natural gradient
      $\delta\theta = -\eta\QFI^{-1}\nabla_\theta\mathcal{L}$
      respects the Fubini--Study metric and converges faster
      on stiff, anisotropic loss landscapes.
    \end{block}
    \begin{alertblock}{Barren plateau warning}
      Adding more layers improves expressivity but risks
      entering the barren plateau regime.
      Physics-inspired ans\"atze mitigate this: use
      $L \leq \log n$ layers with structured entanglement.
    \end{alertblock}
  \end{columns}
\end{frame}

%% ============================================================
\section{Extended Applications}
%% ============================================================

%% ─── Expressivity and Fourier analysis ──────────────────────
\begin{frame}{Expressivity of QPINNs: Fourier Analysis}
  \begin{columns}[T]
  \column{0.52\textwidth}
    \textbf{Theorem (Schuld et al.\ 2021).}
    The output of any QNN with data-encoding gates
    $e^{-ixP/2}$ is a \emph{partial Fourier series}:
    \[
      u_\theta(x)
      = \sum_{\omega\in\Omega}
        c_\omega(\theta)\,e^{i\omega x}
    \]
    where $\Omega\subseteq\{\omega_1{+}\cdots{+}\omega_n:
      \omega_i\in\{0,\pm\tfrac12\}\}$.
    Frequencies fixed by encoding; $c_\omega$ tuned by $\theta$.

    \textbf{For PDEs:} highest frequency in $\Omega$
    limits the shortest spatial scale resolved.

  \column{0.44\textwidth}
    \begin{block}{Error decomposition}
      \[
        \|u_{\theta^\star}-u^\star\|_{L^2}
        \leq
        \varepsilon_{\rm approx}
       +\varepsilon_{\rm opt}
       +\varepsilon_{\rm stat}
      \]
      Expressivity: $\varepsilon_{\rm approx}
       \sim \sum_{\omega\notin\Omega}|\hat{u}^\star_\omega|$;
      sampling: $\varepsilon_{\rm stat}\sim N_r^{-1/2}$.
    \end{block}
    \begin{alertblock}{Circuit sizing}
      For target frequency $\omega^\star$, choose
      $n$ qubits and $L$ re-uploading layers such that
      $Ln \geq 2\omega^\star$.
    \end{alertblock}
  \end{columns}
\end{frame}

%% ─── Time-dependent PDEs ────────────────────────────────────
\begin{frame}{Extension: Time-Dependent PDEs}
  \begin{columns}[T]
  \column{0.52\textwidth}
    For $u(x,t)$, encode space and time jointly:
    \[
      S(x,t) = R_y(\pi x)^{\otimes n_x}
               \otimes R_y(\pi t)^{\otimes n_t}.
    \]
    \textbf{1D heat equation:}
    $\partial_t u - \kappa\partial_{xx}u = 0$

    Input shifts for PDE residual:
    \begin{align*}
      \partial_t u
        &\approx \tfrac12[u(x,t{+}\tfrac\pi2)
                         -u(x,t{-}\tfrac\pi2)]\\
      \partial_{xx}u
        &\approx u(x{+}\tfrac\pi2,t)
                 - 2u(x,t)
                 + u(x{-}\tfrac\pi2,t)
    \end{align*}
    Each residual point: 5 circuit evaluations.

    Add IC loss: $\mathcal{L}_{\rm IC}
    = N_x^{-1}\sum_k|u_\theta(x_k,0)-u_0(x_k)|^2$.

  \column{0.44\textwidth}
    \begin{alertblock}{Dimensional scaling}
      $d$-dimensional PDE: $n_x + n_t = n$ total qubits.
      Classical PINN cost: $O(\varepsilon^{-d})$.
      QPINN cost: $O(d)$ qubits if solution has
      tensor-product structure.
      Advantage largest when $d\gg1$.
    \end{alertblock}
    \begin{block}{Exact: $u(x,t) = e^{-\pi^2t}\sin(\pi x)$}
      Compare QPINN to exact solution; measure
      $L^\infty$ error across the space-time domain.
    \end{block}
  \end{columns}
\end{frame}

%% ─── Schrodinger equation ────────────────────────────────────
\begin{frame}{The Schr\"odinger Equation as a QPINN}
  \begin{columns}[T]
  \column{0.52\textwidth}
    Time-dependent Schr\"odinger equation:
    \[
      i\hbar\partial_t\psi = H\psi,\quad
      H = -\frac{\hbar^2}{2m}\partial_{xx} + V(x).
    \]
    Split $\psi = \psi_R + i\psi_I$:\vspace{-3pt}
    \begin{align*}
      \partial_t\psi_R
        &= \tfrac\hbar{2m}\partial_{xx}\psi_I - V\psi_I/\hbar\\
      \partial_t\psi_I
        &= -\tfrac\hbar{2m}\partial_{xx}\psi_R + V\psi_R/\hbar
    \end{align*}
    Two QPINNs: $\psi_R\!\approx\!\langle Z_0\rangle_{x,t,\theta_R}$,
    $\psi_I\!\approx\!\langle Z_0\rangle_{x,t,\theta_I}$.
    All derivatives by input-shift; $V(x)$ analytically.
  \column{0.44\textwidth}
    \begin{alertblock}{Bridge to Hamiltonian simulation}
      The QPINN Schr\"odinger ansatz is the
      \emph{classical simulation} of quantum evolution.
      On hardware, $e^{-iHt/\hbar}$ is exact --- the QPINN
      bridges PDE solving and digital quantum simulation.
    \end{alertblock}
    \begin{block}{Time-independent: VQE connection}
      $H\psi=E\psi$ solved by minimising
      $E[\psi_\theta]=\langle H\rangle_{\psi_\theta}\ge E_0$
      (Rayleigh quotient = VQE objective).
      QFI/TDVP (Section~3) govern both VQE and QPINN.
    \end{block}
  \end{columns}
\end{frame}

%% ─── Barren plateaus in QPINNs ───────────────────────────────
\begin{frame}{Barren Plateaus in QPINNs: Additional Sources}
  QPINNs inherit the general barren plateau problem
  (Section~4) and have additional sources specific to PDE training:
  \begin{columns}[T]
  \column{0.52\textwidth}
    \textbf{General BP} (random deep circuits):
    \[
      \Variance\!\left(
        \frac{\partial\mathcal{L}}{\partial\theta_\mu}
      \right)\sim\frac{1}{2^n}.
    \]
    \textbf{QPINN-specific sources:}
    \begin{itemize}
      \item Residual $r_k\approx 0$ early in training:
            gradient is small even without BPs
      \item Loss imbalance:
            $\mathcal{L}_r\gg\mathcal{L}_b$
            makes BC enforcement slow
      \item Stiff PDEs: widely varying eigenvalues
            of the differential operator
    \end{itemize}
  \column{0.44\textwidth}
    \begin{block}{QPINN-specific mitigation}
      \begin{enumerate}\itemsep2pt\small
        \item \textbf{Physics ansatz:} mirror Green's function structure
        \item \textbf{Incremental training:} low modes first
        \item \textbf{Adaptive collocation:} focus on large-residual points
        \item \textbf{Loss weighting:} normalise $\mathcal{L}_r$,
              $\mathcal{L}_b$ by gradient norms
        \item \textbf{Natural gradient:}
              $\dot\theta = F^{-1}\nabla_\theta\mathcal{L}$
      \end{enumerate}
    \end{block}
  \end{columns}
\end{frame}

%% ============================================================
\section{Summary, Exercises, and References}
%% ============================================================

%% ─── Comparison table ────────────────────────────────────────
\begin{frame}{QPINN vs.\ Classical PINN: Summary}
\footnotesize
  \begin{center}
  \renewcommand{\arraystretch}{1.0}
  \begin{tabular}{p{2.4cm}p{3.9cm}p{3.9cm}}
    \toprule
    & \textbf{Classical PINN} & \textbf{QPINN} \\
    \midrule
    Ansatz
      & Deep MLP $u_\theta(x)$
      & Circuit $\langle O\rangle_{x,\theta}$ \\
    Spatial deriv.
      & Automatic differentiation
      & Input-shift rule (exact) \\
    Parameter grad.
      & Backpropagation
      & Parameter-shift rule (exact) \\
    Frequency
      & Arbitrary (spectral bias)
      & Bounded, explicit $\Omega$ \\
    Memory
      & $O(N_p)$ params
      & $O(n)$ params, $2^n$ amplitudes \\
    High-$d$ PDE
      & Exponential cost
      & Potential $O(d)$ qubit advantage \\
    Optimiser geom.
      & Euclidean
      & QFI / Fubini--Study (natural grad.) \\
    Advantage
      & Production-ready
      & High-$d$, quantum PDEs \\
    \bottomrule
  \end{tabular}
  \end{center}
  \begin{alertblock}{The central message}\footnotesize
    QPINN = physics-informed QNN: circuit $\approx$ PDE solution;
    spatial derivatives via input-shift rule;
    training minimises residual with exact quantum gradients;
    QFI/TDVP give the natural optimisation geometry.
  \end{alertblock}
\end{frame}

%% ─── Key equations ───────────────────────────────────────────
\begin{frame}{Key Equations}
  \small
  \begin{center}
  \begin{tabular}{ll}
    \toprule
    \textbf{Concept} & \textbf{Formula} \\
    \midrule
    Parameter-shift rule &
      $\dfrac{\partial f}{\partial\theta_j}
       = \dfrac{1}{2}[f(\theta_j{+}\tfrac\pi2)
                     -f(\theta_j{-}\tfrac\pi2)]$ \\[7pt]
    Input-shift rule &
      $\dfrac{\partial u_\theta}{\partial x}
       = \dfrac{1}{2}[u_\theta(x{+}\tfrac\pi2)
                     -u_\theta(x{-}\tfrac\pi2)]$ \\[7pt]
    Laplacian (2nd order) &
      $\partial_{xx}u_\theta
       = u_\theta(x{+}\tfrac\pi2)
         - 2u_\theta(x)
         + u_\theta(x{-}\tfrac\pi2)$ \\[7pt]
    QFI (pure state) &
      $F_{jk} = 4\,\mathrm{Re}[
        \langle\partial_j\psi|\partial_k\psi\rangle
       -\langle\partial_j\psi|\psi\rangle
        \langle\psi|\partial_k\psi\rangle]$ \\[7pt]
    Natural gradient &
      $\btheta\leftarrow\btheta-\eta\,\QFI^{-1}\nabla C$ \\[7pt]
    TDVP &
      $\sum_k F_{jk}\dot{\theta}_k
       = 2\,\mathrm{Im}\langle\partial_j\psi|H|\psi\rangle$ \\[7pt]
    Barren plateau &
      $\Variance[\partial C/\partial\theta_j]\sim 1/2^n$ \\[7pt]
    QPINN loss &
      $\mathcal{L} = \|\mathcal{D}[u_\theta]-f\|^2
                    + \lambda\|u_\theta-g\|^2_{\partial\Omega}$ \\
    \bottomrule
  \end{tabular}
  \end{center}
\end{frame}

%% ─── Exercises ───────────────────────────────────────────────
\begin{frame}{Exercises}
  \begin{enumerate}\itemsep5pt
    \item \textbf{Parameter-shift by hand.}
          For $f(\theta)=\langle 0|R_y(\theta)^\dagger Z R_y(\theta)|0\rangle$,
          apply the parameter-shift rule and verify by direct
          differentiation.

    \item \textbf{QFI of a single qubit.}
          For $|\psi(\theta)\rangle = R_y(\theta)|0\rangle$,
          compute $F_{11}$ via the Fubini--Study formula and
          interpret geometrically.

    \item \textbf{Barren plateau scaling.}
          Numerically estimate
          $\Variance[\partial C/\partial\theta_j]$ for
          $n=2,\ldots,8$ qubit random circuits of depth $d=2n$.
          Verify the $O(1/2^n)$ scaling.

    \item \textbf{Harmonic oscillator QPINN.}
          Modify the Poisson code to solve $u''+\omega^2 u=0$,
          $u(0)=0$, $u'(0)=1$.
          Compare to $u(x)=\sin(\omega x)/\omega$.

    \item \textbf{Second-order input-shift identity.}
          Prove that for $S(x)=e^{-ixP/2}$ with
          eigenvalues $\pm\tfrac12$:
          \[
            \frac{d^2u_\theta}{dx^2}
            = u_\theta(x{+}\tfrac\pi2)
              - 2u_\theta(x)
              + u_\theta(x{-}\tfrac\pi2).
          \]

    \item \textbf{Natural gradient QPINN.}
          Replace Adam with the diagonal QFI natural gradient
          in the Poisson solver and compare convergence.
  \end{enumerate}
\end{frame}

%% ─── References ──────────────────────────────────────────────
\begin{frame}[allowframebreaks]{References}
  \begin{columns}[T]
  \column{0.50\textwidth}
    \textbf{QNN foundations}
    \begin{enumerate}\scriptsize
      \item J.R.\ McClean et al., \emph{Barren plateaus},
            Nat.\ Commun.\ \textbf{9}, 4812 (2018).
      \item M.\ Cerezo et al., \emph{Cost-function-dependent BPs},
            Nat.\ Commun.\ \textbf{12}, 1791 (2021).
      \item J.\ Stokes et al., \emph{Quantum natural gradient},
            Quantum \textbf{4}, 269 (2020).
      \item A.\ P\'erez-Salinas et al., \emph{Data re-uploading},
            Quantum \textbf{4}, 226 (2020).
      \item A.\ Abbas et al., \emph{Power of QNNs},
            Nat.\ Comput.\ Sci.\ \textbf{1}, 403 (2021).
      \item A.\ Havl\'{\i}\v{c}ek et al., \emph{Supervised learning
            with quantum-enhanced feature spaces},
            Nature \textbf{567}, 209 (2019).
      \item S.\ Wang et al., \emph{Noise-induced BPs},
            Nat.\ Commun.\ \textbf{12}, 6961 (2021).
      \item S.\ Sorella, \emph{Stochastic reconfiguration},
            PRL \textbf{80}, 4558 (1998).
    \end{enumerate}
  \column{0.46\textwidth}
    \textbf{Classical PINNs}
    \begin{enumerate}\scriptsize\setcounter{enumi}{8}
      \item M.\ Raissi et al., \emph{Physics-informed NNs},
            J.\ Comput.\ Phys.\ \textbf{378}, 686 (2019).
      \item S.\ Wang et al., \emph{When and why PINNs fail},
            J.\ Comput.\ Phys.\ \textbf{449}, 110768 (2022).
    \end{enumerate}
    \textbf{Quantum PINNs and expressivity}
    \begin{enumerate}\scriptsize\setcounter{enumi}{10}
      \item A.\ Kyriienko et al.,
            \emph{Solving nonlinear DEs with differentiable
            quantum circuits},
            PRA \textbf{103}, 052416 (2021).
      \item M.\ Schuld et al.,
            \emph{Effect of data encoding on expressive power},
            PRA \textbf{103}, 032430 (2021).
      \item M.\ Cerezo et al.,
            \emph{Variational quantum algorithms},
            Nat.\ Rev.\ Phys.\ \textbf{3}, 625 (2021).
    \end{enumerate}
  \end{columns}
\end{frame}

\end{document}
